\def\EditionMode{1}
\def\EditionLabel{Local math-prose proof}
\def\EditionStatus{Proposed wrappers; canonical source unchanged}
\documentclass[10pt,oneside]{book}
\usepackage{fontspec}
\setmainfont{Latin Modern Roman}
\setsansfont{Latin Modern Sans}
\usepackage{amsmath,amssymb,mathtools}
\usepackage{unicode-math}
\setmathfont{Latin Modern Math}
\usepackage{xcolor,tikz-cd,lua-ul,marginnote,pdflscape}
\usetikzlibrary{arrows.meta,decorations.pathmorphing,calc,shapes.geometric}
\tikzcdset{ampersand replacement=\&}
\tikzcdset{every label/.append style={text=black}}
\usepackage[a4paper,left=18mm,right=45mm,top=18mm,bottom=22mm,marginparwidth=28mm,marginparsep=7mm]{geometry}
\usepackage[unicode,colorlinks=true,urlcolor=black,linkcolor=black]{hyperref}
\usepackage{bookmark}
\setlength{\parindent}{0pt}
\setlength{\parskip}{4pt plus 1pt minus 1pt}
\setlength{\emergencystretch}{1.5em}
\raggedbottom
% Include final-line descenders inside the declared text-height rectangle.
\maxdepth=0pt
\definecolor{PSHand}{RGB}{25,58,132}
\definecolor{PSDelete}{RGB}{146,37,37}
\definecolor{PSTypedOld}{RGB}{95,95,95}
\ifnum\EditionMode=2\colorlet{blue}{black}\else\colorlet{blue}{PSHand}\fi
\RenewDocumentCommand{\textcircled}{m}{\tikz[baseline=(PSCircle.base)]\node[draw,circle,inner sep=.5pt,line width=.4pt](PSCircle){#1};}
\newif\ifPSInApparatus
\newcounter{psnote}
\newcounter{psvariant}
\newcounter{pslinkmark}
\newif\ifPSInLink
\newsavebox{\PSNoteBox}
\newdimen\PSNoteLimit
{\footnotesize\global\PSNoteLimit=3\baselineskip}
\newcommand{\CurrentScanURL}{}
\newcommand{\CurrentKey}{}
\newcommand{\CurrentScan}{}
\newcommand{\pssource}[3]{\psflushmargins\clearpage\gdef\CurrentKey{#1}\gdef\CurrentScan{#2}\gdef\CurrentScanURL{#3}\setcounter{psvariant}{0}\typeout{PS-SOURCE: #1}}
\newcommand{\psdivision}[2]{\pdfbookmark[0]{#1}{#2}\par\noindent{\large\bfseries #1\par}\medskip}
\newcommand{\pslinkmarker}[2]{\stepcounter{pslinkmark}\latelua{texio.write_nl('log','PS-LINK-SHIP|#1|\CurrentKey|\arabic{pslinkmark}|#2')}}
\newcommand{\pslink}[2]{\ifPSInLink\pslinkmarker{#1}{nested}#2\else\begingroup\PSInLinktrue\href{\CurrentScanURL}{\pslinkmarker{#1}{outer}#2}\endgroup\fi}
\newcommand{\scanpage}[1]{\ifnum\EditionMode=2\leavevmode\marginpar{\scriptsize\sffamily\pslink{scanpage}{scan #1}}\else\par\noindent{\small\sffamily\pslink{scanpage}{[scan #1]}}\par\medskip\fi}
\newcommand{\mspage}[1]{\ifnum\EditionMode=2\leavevmode\marginpar{\scriptsize\sffamily p.\,#1}\else{\small\sffamily[p.\,#1]}\fi}
\newcommand{\hand}[1]{{\ifnum\EditionMode=2\color{black}\else\color{PSHand}\fi #1}}
\let\over\relax
\NewDocumentCommand{\over}{O{} m m}{{\textcolor{PSTypedOld}{#2}\ensuremath{^{\!\nearrow}}\if\relax\detokenize{#1}\relax\hand{#3}\else #3\fi}}
\NewDocumentCommand{\inserted}{O{} m}{{\if\relax\detokenize{#1}\relax\hand{\ensuremath{\langle}#2\ensuremath{\rangle}}\else\ensuremath{\langle}#2\ensuremath{\rangle}\fi}}
\newcommand{\deleted}[1]{{\color{PSDelete}\ifmmode\text{[del.: }#1\text{]}\else[\textsf{del.:} #1]\fi}}
\newcommand{\xout}[1]{{\color{PSTypedOld}\ifmmode\text{[x: }#1\text{]}\else[\textsf{x:} #1]\fi}}
\let\PSMathUnderline\underline
\RenewDocumentCommand{\underline}{+m}{\ifmmode\PSMathUnderline{#1}\else\underLine{#1}\fi}
\NewDocumentCommand{\uline}{+m}{\underline{#1}}
\NewDocumentCommand{\uuline}{+m}{\ifmmode\PSMathUnderline{\PSMathUnderline{#1}}\else\underLine{\underLine{#1}}\fi}
\newunderlinetype\PSWaveLine{\leaders\hbox{\lower.7ex\hbox{\tikz[x=1pt,y=1pt,baseline=0pt]{\draw[line width=.3pt] (0,0) .. controls (1,1) and (2,1) .. (3,0) .. controls (4,-1) and (5,-1) .. (6,0);}}}}
\NewDocumentCommand{\uwave}{+m}{{\PSWaveLine #1}}
\newcommand{\unsure}[1]{\pslink{unsure}{[#1\ifmmode\,?\else\,?\fi]}}
\newcommand{\illegible}[1]{\ifmmode\text{\pslink{illegible}{[illeg.\ #1]}}\else\pslink{illegible}{[illeg.\ #1]}\fi}
\NewDocumentCommand{\diagram}{+m}{\ifmmode\vcenter{\hbox{#1}}\else\par\medskip{\centering #1\par}\medskip\psflushmargins\fi}
\newenvironment{diagramblock}{\par\begin{center}}{\end{center}\par}
\newcommand{\PSMathNotes}{}
\makeatletter
\NewDocumentCommand{\psplacednote}{m +m}{%
      \sbox{\PSNoteBox}{\parbox{\marginparwidth}{\footnotesize\raggedright\hand{\textsf{m#1.} #2}}}%
      \ifdim\dimexpr\ht\PSNoteBox+\dp\PSNoteBox\relax<\PSNoteLimit
        \leavevmode\marginpar{\usebox{\PSNoteBox}}%
      \else
        \par\nobreak\smallskip\hrule\nobreak\smallskip
        {\footnotesize\textsf{Marginal note m#1.}\quad\hand{#2}\par}
        \nobreak\smallskip\hrule\smallskip
      \fi}
\newcommand{\psflushmargins}{\ifmmode\else\ifinner\else\ifx\PSMathNotes\@empty\else\let\PSTempNotes\PSMathNotes\global\let\PSMathNotes\@empty\PSTempNotes\fi\fi\fi}
\NewDocumentCommand{\margin}{+m}{%
  \ifPSInApparatus\ifmmode\hand{\text{[marg.: }#1\text{]}}\else\hand{\textsf{[marg.: }#1\textsf{]}}\fi
  \else\stepcounter{psnote}%
    \ifmmode
      \text{\textsuperscript{m\arabic{psnote}}}%
      \begingroup\edef\PSQueue{\endgroup\noexpand\g@addto@macro\noexpand\PSMathNotes{\noexpand\psplacednote{\arabic{psnote}}{\unexpanded{$#1$}}}}\PSQueue
    \else
      \textsuperscript{m\arabic{psnote}}\psplacednote{\arabic{psnote}}{#1}%
    \fi
  \fi}
\AtBeginDocument{\everymath\expandafter{\the\everymath\aftergroup\psflushmargins}\everydisplay\expandafter{\the\everydisplay\aftergroup\psflushmargins}}
\makeatother
\newcommand{\psnumber}[2]{\expandafter\gdef\csname psnum@#1\endcsname{#2}}
\NewDocumentCommand{\var}{+m +m m m}{\ifnum\EditionMode=0 #1\else\ifnum\EditionMode=1\ifmmode\text{\textsuperscript{\tiny\csname psnum@#4\endcsname}}\else\textsuperscript{\tiny\csname psnum@#4\endcsname}\fi\fi #2\fi}
\newcommand{\psapparatusstart}{\par\bigskip\hrule\smallskip\PSInApparatustrue\begingroup\footnotesize\raggedright\textsf{Apparatus for scan \CurrentScan}\par}
\newcommand{\psapparatusend}{\par\endgroup\PSInApparatusfalse}
\NewDocumentCommand{\psentry}{m m m +m +m m}{\par\noindent\textbf{#2. #3.}\quad\textit{Typescript:} #4\quad\textit{Corrected:} \if\relax\detokenize{#5}\relax\textnormal{[omitted]}\else #5\fi\quad #6\par}
\newcommand{\psfrontmatter}{%
\hypersetup{pageanchor=false}\begin{titlepage}\raggedright
{\Large Alexander Grothendieck\par}\vspace{12mm}
{\Huge Pursuing Stacks\par}\medskip{\LARGE 1983\par}\vspace{12mm}
{\Large\EditionLabel\par}\vspace{8mm}
{\large\EditionStatus\par}\vfill
Université de Montpellier, Grothendieck archive, cotes 134-2--134-8.\par
Including the 26 source pages of R. Brown's translation at cote 134-2, sites 115--143.\par
\vspace*{2mm}\end{titlepage}\hypersetup{pageanchor=true}
\chapter*{Editorial statement}
AI-transcribed by Claude Opus 5.5 (Anthropic, effort high); checked against the scans, and this edition built, by GPT-6 Astra (OpenAI, ChatGPT Pro mode); prompts and records by Claude (Anthropic).

The witness consulted and not reproduced is the LaTeX edition of Mateo Carmona with the collaboration of Ulrik Buchholtz (arXiv 2111.01000 v2, 2021), extended from the 2015 edition by ``the scrivener'', CC0 1.0.

The translation in the appendix is R. Brown's, prepared with assistance from L. Breen and J-L. Loday, as his note says; the rights in it belong to its translator or his heirs.

The rights in the original work belong to its author or the author's representatives or heirs. Where we hold rights, in the AI-produced transcription, translation, corrections and typesetting, they are dedicated to the public domain under CC0 1.0. This edition is published for educational, non-profit purposes, and we claim no rights we do not hold. If you hold rights in this work and object to its publication here, please open a GitHub issue and we will take it down.

\section*{How to read this form}
\ifnum\EditionMode=0
The diplomatic form preserves typed and handwritten wording, cancellations, insertions, overwritten readings, underlining, doubts, and original page numbers. Blue marks handwriting; red labels handwritten cancellations; gray labels typewriter cancellations and underlying overwritten text. Angle brackets mark insertions; a superscript arrow separates an overwritten reading from its replacement. Typed insertions and typed overwriting remain plain. The disposition of long marginal notes is identified explicitly.
\else
The corrected text applies the author's insertions, replacements, cancellations, and recorded transpositions. It retains his English and notation, including consistent forms such as ``standart'', ``developped'', ``it's'' for ``its'', and ``A la''. An isolated slip is corrected only when specifically recorded and justified. Mathematical anomalies are preserved and discussed in the guide.
\ifnum\EditionMode=1
Blue continues to identify handwriting. A numbered call identifies each maximal editorial operation, with its full diplomatic and corrected readings and reason in the apparatus following the source page. Internal layers of a wholly replaced or deleted passage are preserved within that entry rather than duplicated as separate entries. These entries are generated from the same variants that govern the running text.
\else
This reading form omits the apparatus, cancelled material, superseded readings, insertion brackets, and color distinctions. It retains the resulting text, marginal notes, underlining, and unresolved doubts. Scan and source-page references appear in the margin.
\fi\fi

Every scan mark and remaining doubt links to its source page in the archive viewer. Doubts are editorial reading marks, not emendations. The source-page breaks identify the archival sequence; a source page may occupy more than one printed page. A partial working edition is explicitly labelled and contains only reviewed source pages.
\clearpage}

\begin{document}
\pssource{ps_134-6_061}{134-6/61}{https://grothendieck.umontpellier.fr/web/viewer.html?file=../134-6.pdf\#page=62}
\psnumber{ps_134-6_061-v0001}{1}
\psnumber{ps_134-6_061-v0002}{2}
\psnumber{ps_134-6_061-v0003}{3}
\psnumber{ps_134-6_061-v0004}{4}
\psnumber{ps_134-6_061-v0005}{5}
\psnumber{ps_134-6_061-v0006}{6}
\psnumber{ps_134-6_061-v0007}{7}
\psnumber{ps_134-6_061-v0008}{8}
\psnumber{ps_134-6_061-v0009}{9}
\psnumber{ps_134-6_061-v0010}{10}
\psnumber{ps_134-6_061-v0011}{11}
\psnumber{ps_134-6_061-v0012}{12}
\psnumber{ps_134-6_061-v0013}{13}
\psnumber{ps_134-6_061-v0014}{14}
\psnumber{ps_134-6_061-v0015}{15}
\psnumber{ps_134-6_061-v0016}{16}
\psnumber{ps_134-6_061-v0017}{17}
\psnumber{ps_134-6_061-v0018}{18}
% Pursuing Stacks (cotes 134-2 to 134-8): archive cote 134-6, site page 61 (PDF page 62), typescript page 379.
% typescript. Draft compared: ps5.tex:4443-4519.
% AI-transcribed by Claude Opus 5.5 (claude-opus-5-5, adaptive thinking, effort high), Message Batches API, batch ps3 (msgbatch_01XdFmuVVBAcgc7EFxdtCqsj), 2026-11-12; reviewed against both supplied scan halves; see CHECK_DECISIONS.tsv.
\scanpage{134-6/61}
\mspage{379}

to equivalence)  each left hand factor to the category of functors
from the other functor to (Ab)  which commute with small direct limits
\var{\xout{\unsure{(}}}{}{typewriter deletion}{ps_134-6_061-v0001}( much in the same way as the corresponding relationship between $A^{\wedge}$and
$B^{\wedge}$, with (A\var{\over[typed]{\illegible{1}}{b}}{b}{typewriter overwriting}{ps_134-6_061-v0002})  being replaced by (Sets))\var{\xout{\unsure{)}}}{}{typewriter deletion}{ps_134-6_061-v0003} ,  and that accordingly the
functor (5) is equally fully faithful,  and extends to a fully faithful
functor from $B^{\wedge}_{ab}$ to  \uline{Homadd}$(A^{\wedge}_{ab},(Ab))$ ,  inducing in fact an equivalence
between $B^{\wedge}_{ab}$  and the full subcategory  \hand{\uline{\textcolor{black}{Hom}}}$_{\hand{!}}(A^{\wedge}_{ab},(Ab))$ of \hand{\uline{\textcolor{black}{Hom}}}$(A^{\wedge}_{ab},(Ab))$
made up by all functors  $A^{\wedge}_{ab}\hand{\to}(Ab)$  which commute to small direct limits\hand{.}
But for what we have in mind at present,  th\var{\over[typed]{i}{e}}{e}{typewriter overwriting}{ps_134-6_061-v0004}se niceties are not too re-%
levant yet it seems  -  all what matters is that an object L of Addinf(B)
defines an\var{\xout{functorx}}{}{typewriter deletion}{ps_134-6_061-v0005} exact functor
\[ \overset{\hand{\sim}}{L} \ : \ A^{\wedge}_{ab} \hand{\longrightarrow} (Ab) \ , \]
depending functorially on L,  in an infinitely additive way.  Thus,  as
noted in section 92  (but w\var{\over{ith}{here}}{\hand{here}}{writing over}{ps_134-6_061-v0006} Addinf(B) \var{\deleted{being}}{}{deletion}{ps_134-6_061-v0007}\var{\inserted{was}}{\hand{was}}{insertion}{ps_134-6_061-v0008} replaced by the smaller
category Add(B),  which has turned out unsufficient for our purposes),
whenever we have a chain complex L.  in Addinf(B),  we get a corresponding
functor
\[ \text{(7)} \qquad \overset{\hand{\sim}}{L}. \ : \ A^{\wedge}_{ab} \hand{\longrightarrow} Ch.(Ab) \]
from $A^{\wedge}_{ab}$ to the category of chain complexes of (Ab),  which is moreover
an \hand{\uline{\textcolor{black}{exact}}} functor.  Generalizing slightly the \var{\xout{xxx}}{}{typewriter deletion}{ps_134-6_061-v0009} terminology introduced
in \var{\inserted{section}}{\hand{section}}{insertion}{ps_134-6_061-v0010} 93,  \var{\xout{xxxxx}}{}{typewriter deletion}{ps_134-6_061-v0011} where we restricted to chain complexes with components in
Add(B) rather than in Addinf(B),  we'll say that  $\overset{\hand{\sim}}{L}.$  is an \hand{\uline{\textcolor{black}{abelianizator}}}
\hand{\uline{\textcolor{black}{for}}}  A  ,  if the  following diagram commutes up to \var{\over[typed]{a}{i}}{i}{typewriter overwriting}{ps_134-6_061-v0012}somorphism :
\diagram{
\[ \text{(8)} \qquad
% Hand-drawn arrows are marked blue; typed objects remain plain.
\begin{tikzcd}
A^{\var{\over[typed]{\illegible{1}}{\wedge}}{\wedge}{typewriter overwriting}{ps_134-6_061-v0013}} \arrow[color=blue,r] \arrow[color=blue,d, "\begin{array}{c} Wh_A \\ \hand{\text{(abelianization)}} \end{array}"'] \& Hot_A \arrow[color=blue,r] \& (Hot) \arrow[color=blue,d, "\begin{array}{c} \hand{\text{``absolute''}} \\ \text{abelianization} \\ \text{functor} \end{array}"] \\
A^{\wedge}_{ab} \arrow[color=blue,r, "\hand{\widetilde{L.}}"] \& Ch.(Ab) \arrow[color=blue,r] \& (Hotab)
\end{tikzcd}
\qquad . \]
}
More accurately,  an abelianizator is a pair (L. ,\hand{$\lambda$} ) ,  where \hand{$\lambda$} is an
isomorphism of functors   $A^{\wedge} \hand{\longrightarrow} (Hotab)$  making the diagram commute.
Here \var{\xout{xxx}}{}{typewriter deletion}{ps_134-6_061-v0014} I like to view the abelianization functor
\[ \text{(9)} \qquad (Hot) \hand{\longrightarrow} (Hotab) \overset{def}{=} D.(Ab) \]
as the one described directly in section (100) via integrators of arbi-%
trary modelizing objects i\var{\xout{h}}{}{typewriter deletion}{ps_134-6_061-v0015}n (Cat),  without any reference to an auxiliary
test category such as $\Delta$ or the like. \var{\hand{\tikz[x=.4em,y=.4em,baseline=0]{\draw (-2,0)--(0,0)--(0,1.3)--(2,1.3);\draw (.15,0)--(.15,1.3);}}}{}{placement}{ps_134-6_061-v0016}
% The stepped paragraph locator runs below like, then rises to a stroke below without any.

\var{T\(\underset{\text{h}}{\text{h}}\)e}{The}{typed correction}{ps_134-6_061-v0017} point \var{\inserted{of (8)}}{\hand{of (8)}}{insertion}{ps_134-6_061-v0018} is that  via an
"abelianizator" for A,  we want to be able to give  a)  a \hand{\uline{\textcolor{black}{simultaneous}}}
handy expression,  in terms of ``computable'' chain complexes in (Ab),  of

\margin{61}

\psapparatusstart
\psentry{ps_134-6_061-v0001}{1}{typewriter deletion}{\xout{\unsure{(}}}{}{The typewriter strikes out this material; it leaves the corrected running text.}
\psentry{ps_134-6_061-v0002}{2}{typewriter overwriting}{\over[typed]{\illegible{1}}{b}}{b}{The superposed typed reading replaces the underlying reading.}
\psentry{ps_134-6_061-v0003}{3}{typewriter deletion}{\xout{\unsure{)}}}{}{The typewriter strikes out this material; it leaves the corrected running text.}
\psentry{ps_134-6_061-v0004}{4}{typewriter overwriting}{\over[typed]{i}{e}}{e}{The superposed typed reading replaces the underlying reading.}
\psentry{ps_134-6_061-v0005}{5}{typewriter deletion}{\xout{functorx}}{}{The typewriter strikes out this material; it leaves the corrected running text.}
\psentry{ps_134-6_061-v0006}{6}{writing over}{\over{ith}{here}}{\hand{here}}{The superposed handwritten reading replaces the underlying reading.}
\psentry{ps_134-6_061-v0007}{7}{deletion}{\deleted{being}}{}{The author cancels this material by hand; it leaves the corrected running text.}
\psentry{ps_134-6_061-v0008}{8}{insertion}{\inserted{was}}{\hand{was}}{The handwritten insertion is incorporated at its indicated position; insertion brackets leave the running text.}
\psentry{ps_134-6_061-v0009}{9}{typewriter deletion}{\xout{xxx}}{}{The typewriter strikes out this material; it leaves the corrected running text.}
\psentry{ps_134-6_061-v0010}{10}{insertion}{\inserted{section}}{\hand{section}}{The handwritten insertion is incorporated at its indicated position; insertion brackets leave the running text.}
\psentry{ps_134-6_061-v0011}{11}{typewriter deletion}{\xout{xxxxx}}{}{The typewriter strikes out this material; it leaves the corrected running text.}
\psentry{ps_134-6_061-v0012}{12}{typewriter overwriting}{\over[typed]{a}{i}}{i}{The superposed typed reading replaces the underlying reading.}
\psentry{ps_134-6_061-v0013}{13}{typewriter overwriting}{$\over[typed]{\illegible{1}}{\wedge}$}{$\wedge$}{The superposed typed reading replaces the underlying reading.}
\psentry{ps_134-6_061-v0014}{14}{typewriter deletion}{\xout{xxx}}{}{The typewriter strikes out this material; it leaves the corrected running text.}
\psentry{ps_134-6_061-v0015}{15}{typewriter deletion}{\xout{h}}{}{The typewriter strikes out this material; it leaves the corrected running text.}
\psentry{ps_134-6_061-v0016}{16}{placement}{\hand{\tikz[x=.4em,y=.4em,baseline=0]{\draw (-2,0)--(0,0)--(0,1.3)--(2,1.3);\draw (.15,0)--(.15,1.3);}}}{}{The stepped pen mark indicates the paragraph division already present in the reading.}
\psentry{ps_134-6_061-v0017}{17}{typed correction}{T\(\underset{\text{h}}{\text{h}}\)e}{The}{The lower h is a displaced additional impression beneath the intact baseline word.}
\psentry{ps_134-6_061-v0018}{18}{insertion}{\inserted{of (8)}}{\hand{of (8)}}{The handwritten insertion is incorporated at its indicated position; insertion brackets leave the running text.}
\psapparatusend

\end{document}
